\scanpage{69}
\mspage{22}
\uline{Th. de recollement des tubes} (\unsure{1\textsuperscript{ère}|1\textsuperscript{ière}|p\textsuperscript{re}} version) Considérons le diagramme de type $\widetilde I$, formé par les $X_i^*,V_{i,j}^*,V_{i,j}$ ($i,j\in I$, $j$ un successeur de $i$) et les trois types de flèches \unsure{y} figurant (en trait plein) ds (17). Alors $X$ s'identifie au \deleted{topos} \inserted{topos} $\underset{\longrightarrow}{\mathrm{Lim}}$ de ces topos.

\margin{Pour que cet énoncé puisse prendre un sens précis, il \unsure{faut} sans doute supp. $X$ \uline{sobre}, i.e. déterminé par $\operatorname{Top}(X)\ldots$}

\deleted{Faute de disposer de définitions en forme pour les voisinages tubulaires, je vais faire une démonstration, lorsqu'on \deleted{\unsure{pose}} $X$ remplace \deleted{\unsure{les|chaque}} $V_{ij}$ par des \inserted{vrais} vrais ouverts de $X_j$\unsure{,} de la forme $\bar V_{i,j}\setminus R_{i,j}\cap\bar V_{i,j}$, où $\bar V_{i,j}$ est un voisinage de $X_i$ dans $X_j$, et qu'on pose $\unsure{U_{ij}^*|V_{ij}^*}=U_{ij}\setminus X_i^*$. \deleted{Alors} \deleted{des \unsure{voisinages} $V_{ij}$}}

On va en donner une démonstration heuristique ci-dessous. Mais avant de le faire, je vais le généraliser, en un énoncé qui permet d'exprimer \unsure{égal\textsuperscript{t}|égalem\textsuperscript{t}} les ``topos canoniques'' associés à la \unsure{st.tific\textsuperscript{n}|st.tificat\textsuperscript{n}}, en termes du système inductif de type $\widetilde I$ formé des topos élémentaires.

\margin{68}
